\subsection{\tpdf{$\nCat{n}$, monoidal $n$-structures, and categorical dualities}}

Recall that the category $\nCat{1}=\Cat$ is canonically enriched over itself using the functor categories $\uCat(\cC,\cD)$. It also has a Kelly tensor product given by the cartesian product $\cC \times \cD$ of two categories, which forms a monoidal structure. Therefore, it is natural to ask whether this enrichment and the monoidal structure are compatible. Indeed, we can enrich the monoidal structure to a monoidal $2$-structure\footnote{The term $2$-monoidal category can be confused with $2$-tuply monoidal $1$-category which is far different from the $1$-tuply monoidal $2$-category as described in our case\cite{aguiarmon2010,yau2024bimodcats}.} in a canonical way. There exists a $\uCat$ enriched adjunction
\[\uCat(\cC \times \cD, \cE) \cong \uCat(\cC,\uCat(\cD,\cE))\cdpunct{,}\]
by using a general result which holds for any closed symmetric monoidal category category $\ucV$ \cite[Prop.~3.1.11]{yau2024bimodcats}. In this enriched adjunction, both the product and the internal hom functors are $2$-functors. Furthermore, the left unitor, the right unitor, the associator, and the symmetry are now all $2$-natural transformations. This defines, naturally, a symmetric cartesian closed monoidal $2$-structure. 

\begin{lemma}
	The tuple $(\Cat, \times, \{*\}, a_{\Cat},\lambda_{\Cat}, \rho_{\Cat})$ defines a symmetric cartesian closed monoidal structure on $\Cat$. Furthermore, this monoidal structure can be $\Cat$-enriched by canonically extending the monoidal data $\{\times , a_{\Cat} ,\lambda_{\Cat}, \rho_{\Cat}\}$ to a strict $2$-functor and strict $2$-natural transformations. All the coherence relations remain unchanged. This makes $\Cat$ a monoidal $2$-category.
\end{lemma}


As a consequence, it is natural question to ask whether the same result extends to $\nCat{n}$ or $\nCat{\omega}$.  The following general lemma answers this in affirmative. 

\begin{lemma}
	Let $\cV$ be a locally small, complete, and cocomplete cartesian closed monoidal cateogry. Then, the category $\cVCat$ (of small $\cV$-categories) is also locally small, complete, cocomplete, and cartesian closed. Furthermore, the cartesian product of $\cVCat$ is precisely the Kelly tensor product, as given by the cartesian product of sets of objects and the corresponding products of hom-objects.
\end{lemma} 

\begin{proof}
	Given any symmetric closed monoidal category $\cV$, the (co)completeness of $\cV$ implies the (co)cocompleteness of $\cVCat$ by results of \cite{wolff1974,betti1983,kellyvcatpres2001}. And, the Kelly tensor product of $\cVCat$ forms a symmetric closed monoidal structure by \cite[Sec.~2.3]{kellyect}. We are only left to verify that the kelly tensor product is indeed the cartesian product. This last part essentially follows from the monadicity of $\cVCat$ over $\cVGraph$. 
\end{proof}

In particular, this lemma provides the following desired corollary.

\begin{corollary}
	The tuple $(\nCat{n}, \times, \{*\}, a_{\nCat{n}},\lambda_{\nCat{n}}, \rho_{\nCat{n}})$ canonically defines a symmetric cartesian closed monoidal $(n+1)$-structure on $\nCat{n}$.
\end{corollary}

Finally, this result also extends to $\nCat{\omega}$, which can be defined as a form of colimit of the enrichment process. In particular, the category of all (small) $\nCat{\omega}$-enriched category is just the category $\nCat{\omega}$ itself and is cartesian closed. The interested reader maybe look up the details in \cite{leinster2004,streetncat1987}.

Using suitable universes, it is also clear that each monoidal $(n+1)$-category $\nCat{n}$ is equivalently a \enquote{monoidal object} inside the monoidal $(n+2)$-category $\nCat{(n+1)}$. This provides our first example of a \emph{nested algebraic structure}.

Another primary example is given by the duality involution of strict $n$-categories. Consider the $1$-functor $\op_1\colon \nCat{1} \rightarrow \nCat{1}$ which sends each category $\cC$ to its opposite category $\cC^{\op}$. This is a duality involution $1$-functor on $\nCat{1}$, that is, $\op_1^2=1$. However, as $\nCat{1}$ is naturally a strict $2$-category, we may wish to upgrade this duality involution to a strict $2$-functor. This is not possible. What we instead have is a strict $2$-functor $\op_1 \colon \op_2(\nCat{1}) \rightarrow \nCat{1}$, where $\op_2$ is a strict $1$-functor $\op_2 \colon \nCat{2}_{\leq 1} \rightarrow \nCat{2}_{\leq 1}$ and $(-)_{\leq i}$ denotes truncation of categorical structures. Moreover, the duality of $\op_1$ is now expressed by the equation $\op_1 \circ \op_2(\op_1)=1$.

Thus, $\op_1$ duality is fully defined only in the presence of a higher dimensional $\op_2$ duality. More generally, let $\op_{n,k}$ denote the $k$-cell duality on $n$-categories for $1 \leq k \leq n$. All the $1$-functors $(\op_{n,k})_{\leq 1} \colon n\Cat_{\leq 1} \rightarrow n\Cat_{\leq 1}$ are definable readily. See \cite{ara2015autequivinfcat,markakis2025hominfcat} for example. This leads to the following natural proposition.

\begin{proposition} \label{prop:opdualities}
	For each $1 \leq k \leq n$, the $1$-functor $(\op_{n,k})_{\leq 1} \colon n\Cat_{\leq 1} \rightarrow n\Cat_{\leq 1}$ extends to a corresponding $(n+1)$-functor $\op_{n,k} \colon \op_{n+1,k+1}(n\Cat) \rightarrow n\Cat$.
\end{proposition}

Note that both $\nCat{n}$ and $\op_{n+1,k+1}(\nCat{n})$ are enriched over $\nCat{n}$ by definition as they are $(n+1)$-categories. This means that our primary goal is to extend the given $1$-functors to corresponding $\nCat{n}$-enriched functor. Consequently, we first state the following more general lemma inspired by an example in \cite[Ex.~4.6]{shulman2018contenrich}.

\begin{lemma}
	Let $(\cV, \otimes, I)$ be a closed symmetric monoidal category equipped with a strict symmetric monoidal functor $o\colon \cV \rightarrow \cV$ satisfying $o^2 = 1$. Then, there exists a canonical $\cV$-enriched functor $o \colon o_*(\ucV) \rightarrow \ucV$, which extends the $1$-functor $o \colon \cV \rightarrow \cV$.
\end{lemma}

\begin{proof}
	$\cV$ is enriched, tensored, and cotensored over itself because it is a closed symmetric monoidal category. More precisely, $\cV$ extends to a symmetric monoidal $\cV$-enriched category \cite[Thm.~3.3.2]{yau2024bimodcats}. Thus, we can do a change of base of enrichement to produce a $\cV$-enriched category $o_*(\ucV)$. Note that the underlying category of $o_*(\ucV)$ is still $\cV$ itself because $o$ is a strict monoidal functor. 
	
	
	
	The key idea is to note that $o_*(\ucV)$ is also tensored over $\ucV$. We define the tensoring operation using the functor $\triangleleft \colon \cV \times \cV \xrightarrow{o \times 1}  \cV \times \cV \xrightarrow{\otimes} \cV$. Our goal is now to show that there is a $\cV$-isomorphism
	\[o_*(\ucV)(A \triangleleft X, Y) \cong \ucV(A,o_*(\ucV)(X,Y))\]
	natural in all of $A,X,Y$. This is a straightforward exercise using the Yoneda lemma. Thus, we have established that $\triangleleft$ is indeed a valid tensoring operation.
	
	Finally, we note the following commuting diagram
	% https://q.uiver.app/#q=WzAsNyxbMSwwLCJcXGNWIFxcdGltZXMgXFxjViJdLFswLDEsIlxcY1YgXFx0aW1lcyBcXGNWIl0sWzEsMiwiXFxjViJdLFs0LDIsIlxcY1YiXSxbNCwwLCJcXGNWIFxcdGltZXMgXFxjViJdLFs2LDAsIlggXFxvdGltZXMgbyhZKSJdLFs2LDIsIm8oWCBcXHRyaWFuZ2xlbGVmdCBZKSJdLFswLDEsIm8gXFx0aW1lcyAxIiwyXSxbMSwyLCJcXG90aW1lcyIsMl0sWzIsMywibyIsMl0sWzAsNCwiMSBcXHRpbWVzIG8iXSxbNCwzLCJcXG90aW1lcyAiXSxbMCwyLCJcXHRyaWFuZ2xlbGVmdCIsMl0sWzUsNiwiIiwwLHsibGV2ZWwiOjIsInN0eWxlIjp7ImhlYWQiOnsibmFtZSI6Im5vbmUifX19XV0=
	\[\begin{tikzcd}[ampersand replacement=\&,cramped]
		\& {\cV \times \cV} \&\&\& {\cV \times \cV} \&\& {X \otimes o(Y)} \\
		{\cV \times \cV} \\
		\& \cV \&\&\& \cV \&\& {o(X \triangleleft Y)}
		\arrow["{1 \times o}", from=1-2, to=1-5]
		\arrow["{o \times 1}"', from=1-2, to=2-1]
		\arrow["\triangleleft"', from=1-2, to=3-2]
		\arrow["{\otimes }", from=1-5, to=3-5]
		\arrow[equals, nfold, from=1-7, to=3-7]
		\arrow["\otimes"', from=2-1, to=3-2]
		\arrow["o"', from=3-2, to=3-5]
	\end{tikzcd}\cdpunct[8pt]{,}\]
	which provides a canonical identity left tensorial strength for the unenriched functor $o$. This implies the existence of a canonical $\cV$-enriched functor $o \colon o_*(\ucV) \rightarrow \ucV$ extending the $1$-functor $o \colon \cV \rightarrow \cV$. For more details, see \cite[Def.~3.1.1 and Prop~3.2.4]{segrt2012strengthsmorita}. Explicitly, the action on hom objects is given by the adjoint to composite 
	% https://q.uiver.app/#q=WzAsMyxbMCwxLCJvXyooXFx1Y1YpKFgsWSkgXFxvdGltZXMgb1giXSxbMiwwLCJvKG9fKihcXHVjVikoWCxZKSBcXHRyaWFuZ2xlbGVmdCBYKSJdLFs0LDEsIm9ZIl0sWzAsMSwiIiwwLHsibGV2ZWwiOjIsInN0eWxlIjp7ImhlYWQiOnsibmFtZSI6Im5vbmUifX19XSxbMSwyLCJvKF57XFx0cmlhbmdsZWxlZnR9XFxlcHNpbG9uXllfWCkiXSxbMCwyLCJvKF57XFxvdGltZXN9XFxlcHNpbG9uXllfWCkiLDJdXQ==
	\[\begin{tikzcd}[ampersand replacement=\&,cramped]
		\&\& {o(o_*(\ucV)(X,Y) \triangleleft X)} \&\& \\
		{o_*(\ucV)(X,Y) \otimes oX} \&\&\&\& oY
		\arrow["{o(^{\triangleleft}\epsilon^Y_X)}", from=1-3, to=2-5]
		\arrow[equals, nfold, from=2-1, to=1-3]
		\arrow["{o(^{\otimes}\epsilon^Y_X)}"', from=2-1, to=2-5]
	\end{tikzcd}\cdpunct[8pt]{.}\]
\end{proof}

\begin{comment}
	We can now define the tensor as a $\cV$-enriched functor. 
	
	Note that $o_*(\ucV)$ is also cotensored over $\ucV$ using $- \pitchfork - \colon \cV^{\op} \times \cV \xrightarrow{o^{\op} \times 1} \cV^{\op} \times \cV \xrightarrow{\ucV(-,-)} \cV$. 
\end{comment}

Now, we can use this lemma in our context to obtain the desired result.

\begin{proof}[Proof of \cref{prop:opdualities}]
	This readily follows once we note that $\op_{n+1,k+1}(\nCat{n}) = {(\op_{n,k})}_*(\nCat{n})$ and that all the $1$-functors $(\op_{n,k})_{\leq 1}$ (as defined in \cite[Sec.~5]{markakis2025hominfcat}) are all strictly symmetric and strictly involutive in the desired sense. 
\end{proof}

We can also provide a more concrete description of these $(n+1)$-functors on the hom $n$-categories. Notice that the diagram % https://q.uiver.app/#q=WzAsMixbMCwwLCJbMF0iXSxbMiwwLCJbMV0iXSxbMCwxLCIiLDAseyJvZmZzZXQiOi0yfV0sWzAsMSwiIiwyLHsib2Zmc2V0IjoyfV0sWzEsMF1d
$\begin{tikzcd}[ampersand replacement=\&,cramped]
	{[0]} \&\& {[1]}
	\arrow[shift left=2, from=1-1, to=1-3]
	\arrow[shift right=2, from=1-1, to=1-3]
	\arrow[from=1-3, to=1-1]
\end{tikzcd}$ defines a cocategory object in $\Cat$. In particular, this allows one to compute the functor categories $\cD^{\cC}$ as $(\cD^{\cC})_0 \cong \Cat([0] \times \cC, \cD)$ and $(\cD^{\cC})_1 \cong \Cat([1] \times \cC, \cC) $. All of the source, target, identity, and composition maps for $\cD^{\cC}$ are derived from the cocategory object structure of  % https://q.uiver.app/#q=WzAsMixbMCwwLCJbMF0iXSxbMiwwLCJbMV0iXSxbMCwxLCIiLDAseyJvZmZzZXQiOi0yfV0sWzAsMSwiIiwyLHsib2Zmc2V0IjoyfV0sWzEsMF1d
$\begin{tikzcd}[ampersand replacement=\&,cramped]
	{[0]} \&\& {[1]}
	\arrow[shift left=2, from=1-1, to=1-3]
	\arrow[shift right=2, from=1-1, to=1-3]
	\arrow[from=1-3, to=1-1]
\end{tikzcd}$. 

More generally, the category of all $n$-categories $n\Cat$ has a $n$-cocategory object which can be generated from the diagram $G_0 \rightrightarrows G_1 \rightrightarrows G_2 \cdots \rightrightarrows G_n$, where each $G_i$ is the $i$-globe. This essentially follows from \cite[Prop.~3.14 and Prop.~3.15]{ara2011thetacatcocat} by considering the Globe category embedding $\bG_{n} \rightarrow \Theta_n$.


For each $1\leq k\leq n$, the globe $G_i$ has canonical isomorphisms $\op_{n,k}(G_i) \cong G_i$  induced by the corresponding duality on the $n$-Globe category $\op_{n,k} \colon \bG_n \rightarrow \bG_n$, which swaps the source map and target map that goes into the $k$th level. Indeed, the corresponding $(\op_{n,k})_{\leq 1}$ dualities on $n$-globular presheaves or $n$-categories are also precisely induced from the ones on the Globe category. See \cite[Sec.~5]{markakis2025hominfcat} for more details. Furthermore, like any arbitrary functor, the globe category dualities are compatible with the diagonal map $\bG_n \rightarrow \bG_n \times \bG_n$ and the swap map $\bG_n \times \bG_n \xrightarrow{s} \bG_n \times \bG_n$. Thus, demonstrating the strict symmetric monoidal nature of the induced dualities.

Using this, we can verify that the action of $\op_{n,k}$ on the $i$-transfor  $G_i \times \cC \rightarrow \cD$ is given by
% https://q.uiver.app/#q=WzAsNixbMCwwLCJHX2kgXFx0aW1lcyBDIl0sWzMsMCwiXFxvcF97bixrfUdfaSBcXHRpbWVzIFxcb3Bfe24sa31DIl0sWzUsMCwiR19pIFxcdGltZXMgXFxvcF97bixrfUMiXSxbMCwxLCJEIl0sWzMsMSwiXFxvcF97bixrfUQiXSxbNSwxLCJcXG9wX3tuLGt9RCJdLFswLDNdLFsxLDRdLFsyLDVdLFs2LDcsIihcXG9wX3tuLGt9KV97XFxsZXEgMX0iLDAseyJzaG9ydGVuIjp7InNvdXJjZSI6MjAsInRhcmdldCI6MjB9LCJsZXZlbCI6MSwic3R5bGUiOnsidGFpbCI6eyJuYW1lIjoibWFwcyB0byJ9fX1dLFs3LDgsIlxcY29uZyIsMCx7InNob3J0ZW4iOnsic291cmNlIjoyMCwidGFyZ2V0IjoyMH0sImxldmVsIjoxLCJzdHlsZSI6eyJ0YWlsIjp7Im5hbWUiOiJtYXBzIHRvIn19fV1d
\[\begin{tikzcd}[ampersand replacement=\&,cramped]
	{G_i \times C} \&\&\& {\op_{n,k}G_i \times \op_{n,k}C} \&\& {G_i \times \op_{n,k}C} \\
	D \&\&\& {\op_{n,k}D} \&\& {\op_{n,k}D}
	\arrow[""{name=0, anchor=center, inner sep=0}, from=1-1, to=2-1]
	\arrow[""{name=1, anchor=center, inner sep=0}, from=1-4, to=2-4]
	\arrow[""{name=2, anchor=center, inner sep=0}, from=1-6, to=2-6]
	\arrow["{(\op_{n,k})_{\leq 1}}", between={0.2}{0.8}, maps to, from=0, to=1]
	\arrow["\cong", between={0.2}{0.8}, maps to, from=1, to=2]
\end{tikzcd}\cdpunct[8pt]{.}\]
This ensures that $\op_{n,k}$ changes the direction of only the $k$-transfor. For instance, $\op_{1,1}$ acting on the $1$-functor $[1] \times \cC \rightarrow \cD$ changes only the direction of $[1]$, which is the direction of natural transformations. 

Ultimately, each lower dimensional duality is fully definable only when the higher dimensional duality is defined, at least in a $1$-truncated form. These nested structures can be better formulated and understood using the microcosm principle

\subsection{Mircocosm principle}

The initial statement of the microcosm principle goes back to Baez and Dolan, abstracting a general guiding principle from the observation that monoids are definable in any monoidal category. 

\begin{definition}[Microcosm principle, Baez and Dolan]
	The microcosm principle suggests that, in order to define a generic $(n-1)$-coherent algebraic structure inside an $n$-category, the category itself should be equipped with a categorified $n$-coherent version of the same algebraic structure.
\end{definition} 

A precise formulation of the microcosm principle is given in \cite[Def.~57]{baezhda31998} using the opetopic foundations of $n$-categories. The key idea is to define internal algebraic structures by using maps out of the terminal categorified algebraic structure. However, this result is not immediately applicable to our usual algebraic foundations. Nor does it address whether slicing with respect to an internal algebraic object is always possible. Indeed, consider a monoid object $M$ in a monoidal category $\cV$. The slice category $\cV/M$ admits a canonical monoidal category structure with the sliced tensor product $X \otimes Y \rightarrow M \otimes M \rightarrow M$. This example suggests another general principle governing these slice constructions. Accordingly, we propose the following variant of the microcosm principle.

\begin{definition}[Sliced microcosm principle]
	Let $A$ be a $(n-1)$-coherent algebraic structure inside a $n$-category $\cC$ equipped with the categorified $n$-coherent algebraic structure. The slice $n$-category $\cC/A$ admits a canonical sliced $n$-coherent algebraic structure.
\end{definition}

We emphasize that both of these are merely guiding principles rather than theorems. For instance, applying the microcosm principle on the Lawvere theory of groups doesn't yield a meaningful notion of a group object inside a \enquote{groupoidal category}. The key obstruction for groups is the necessity of the copy operation $\Delta \colon G \rightarrow  G \times G$ and the deletion operation $G \rightarrow 1$ for defining the inverse axioms. 

More generally, the microcosm principle is incompatible with algebraic theories that demands some form of copying or deletion of data. In contrast, it is well suited for theories, such as that of monoids, which require no copying or deletion. We may, in general, expect the microcosm principle to be compatible with any \enquote{linear} multisorted algebraic theory. Furthermore, one should also expect the existence of sliced algebraic structures in such a case.

\subsection{Pseudomonoids and pseudomodules}

Monoidal categories, module categories, and their dualities are important in the latter part of this paper. Therefore, it is useful to introduce more abstract definitions of pseudomonoid objects and pseudomodule objects, both for conceptual clarity and to provide another useful example of the microcosm principle.

Pseudomonoids and pseudomodules are $1$-level higher weak categorification of the idea of a monoid object and a module over a monoid object. These objects can be defined in a general monoidal bicategory $\cV$. We shall view $\cV$ as a one-object tricategory \cite[Def.~4.1]{gurski2013tricatcoh}. Furthermore, all diagrams are read inside the one object cubical tricategory $\st \cV$ (see \cite[Thm.~8.4]{gurski2013tricatcoh} and \cref{appsection:mondiag} for explanation).

\begin{definition}
	A pseudomonoid in a monoidal bicategory $\cV$ is given by the data $(M,\mu, e, a, \lambda, \rho)$ such that
	\begin{enumerate}
		\item $M \in \cV$ is an object, $\mu$ is a $1$-cell $\mu \colon M \otimes M \rightarrow M$, and $e$ is a $1$-cell $e \colon I \rightarrow M$,
		\item $a,\lambda,\rho$ are invertible $2$-cells in $\cV$ which are defined using the following diagrams:
		% https://q.uiver.app/#q=WzAsOSxbMCwwLCJNIFxcb3RpbWVzIChNIFxcb3RpbWVzIE0pIl0sWzAsMiwiKE0gXFxvdGltZXMgTSkgXFxvdGltZXMgTSJdLFswLDQsIk0gXFxvdGltZXMgTSJdLFsyLDQsIk0iXSxbMiwwLCJNIFxcb3RpbWVzIE0iXSxbNCwwLCJNIFxcb3RpbWVzIEkiXSxbNCw0LCJJIFxcb3RpbWVzIE0iXSxbNCwyLCJNIFxcb3RpbWVzIE0iXSxbNiwyLCJNIl0sWzAsNCwiMSBcXG90aW1lcyBcXG11Il0sWzQsMywiXFxtdSJdLFsyLDMsIlxcbXUiLDJdLFsxLDIsIlxcbXUgXFxvdGltZXMgMSIsMl0sWzEsMCwiYV97XFxjVn0iXSxbNSw4LCJcXHJob197XFxjVn0iXSxbNSw3LCIxIFxcb3RpbWVzIGUiLDJdLFs2LDcsImUgXFxvdGltZXMgMSJdLFs2LDgsIlxcbGFtYmRhX3tcXGNWfSIsMl0sWzcsOCwiXFxtdSJdLFsyLDQsImEiLDIseyJzaG9ydGVuIjp7InNvdXJjZSI6MjAsInRhcmdldCI6MjB9LCJsZXZlbCI6Mn1dLFs3LDE0LCJcXHJobyIsMCx7InNob3J0ZW4iOnsic291cmNlIjoyMCwidGFyZ2V0IjoyMH19XSxbNywxNywiXFxsYW1iZGEiLDIseyJzaG9ydGVuIjp7InNvdXJjZSI6MjAsInRhcmdldCI6MjB9fV1d
		\[\begin{tikzcd}
			{M \otimes (M \otimes M)} && {M \otimes M} && {M \otimes I} && \\
			\\
			{(M \otimes M) \otimes M} &&&& {M \otimes M} && M \\
			\\
			{M \otimes M} && M && {I \otimes M}
			\arrow["{1 \otimes \mu}", from=1-1, to=1-3]
			\arrow["\mu", from=1-3, to=5-3]
			\arrow["{1 \otimes e}"', from=1-5, to=3-5]
			\arrow[""{name=0, anchor=center, inner sep=0}, "{\rho_{\cV}}", from=1-5, to=3-7]
			\arrow["{a_{\cV}}", from=3-1, to=1-1]
			\arrow["{\mu \otimes 1}"', from=3-1, to=5-1]
			\arrow["\mu", from=3-5, to=3-7]
			\arrow["a"', between={0.2}{0.8}, Rightarrow, from=5-1, to=1-3]
			\arrow["\mu"', from=5-1, to=5-3]
			\arrow["{e \otimes 1}", from=5-5, to=3-5]
			\arrow[""{name=1, anchor=center, inner sep=0}, "{\lambda_{\cV}}"', from=5-5, to=3-7]
			\arrow["\rho", between={0.2}{0.8}, Rightarrow, from=3-5, to=0]
			\arrow["\lambda"', between={0.2}{0.8}, Rightarrow, from=3-5, to=1]
		\end{tikzcd}\cdpunct[8pt]{,}\]
		\item the associator $a$ and the unitors $\lambda,\rho$ are subject to the pentagon axiom (\cref{app:pentagon}) and the triangle axiom (\cref{app:triangle}).
	\end{enumerate}
\end{definition}

\begin{definition}
	A (lax) $1$-morphism $(f,\chi,\iota) \colon (M,\mu,e,a,\lambda,\rho) \rightarrow (M', \mu',e',a',\lambda',\rho')$ is given by
	\begin{enumerate}
		\item a $1$-cell $f : M \rightarrow M'$,
		\item two $2$-cell $\chi$ and $\iota$ as depicted by the following diagram:
		% https://q.uiver.app/#q=WzAsNyxbMCwwLCJNIFxcb3RpbWVzIE0iXSxbMiwwLCJNJyBcXG90aW1lcyBNJyJdLFswLDIsIk0iXSxbMiwyLCJNJyJdLFs0LDIsIk0iXSxbNiwyLCJNJyJdLFs0LDAsIkkiXSxbMCwxLCJmIFxcb3RpbWVzIGYiXSxbMCwyLCJcXG11IiwyXSxbMSwzLCJcXG11JyJdLFsyLDMsImYiLDJdLFsxLDIsIlxcY2hpIiwyLHsic2hvcnRlbiI6eyJzb3VyY2UiOjIwLCJ0YXJnZXQiOjIwfSwibGV2ZWwiOjJ9XSxbNCw1LCJmIiwyXSxbNiw0LCJlIiwyXSxbNiw1LCJlJyJdLFsxNCw0LCJcXGlvdGEiLDIseyJzaG9ydGVuIjp7InNvdXJjZSI6MjAsInRhcmdldCI6MjB9fV1d
		\[\begin{tikzcd}[ampersand replacement=\&]
			{M \otimes M} \&\& {M' \otimes M'} \&\& I \&\& \\
			\\
			M \&\& {M'} \&\& M \&\& {M'}
			\arrow["{f \otimes f}", from=1-1, to=1-3]
			\arrow["\mu"', from=1-1, to=3-1]
			\arrow["\chi"', between={0.2}{0.8}, Rightarrow, nfold, from=1-3, to=3-1]
			\arrow["{\mu'}", from=1-3, to=3-3]
			\arrow["e"', from=1-5, to=3-5]
			\arrow[""{name=0, anchor=center, inner sep=0}, "{e'}", from=1-5, to=3-7]
			\arrow["f"', from=3-1, to=3-3]
			\arrow["f"', from=3-5, to=3-7]
			\arrow["\iota"', between={0.2}{0.8}, Rightarrow, nfold, from=0, to=3-5]
		\end{tikzcd}\cdpunct[8pt]{,}\]
		\item such that the two $2$-cells satisfy the morphism hexagon (\cref{app:morphhexagon}), morphism left square (\cref{app:morphlsq}), and morphism right square (\cref{app:morphrsq}).
	\end{enumerate}
\end{definition}

\begin{definition}
	A $2$-morphism $r \colon (f,\chi_{f},\iota_{f}) \Rightarrow (g,\chi_{g},\iota_{g})$ is given by
	\begin{enumerate}
		\item a $2$-cell $r \colon f \Rightarrow g$,
		\item such that the following two identities hold:
		% https://q.uiver.app/#q=WzAsOSxbMCwwLCJNIFxcb3RpbWVzIE0iXSxbMiwwLCJNJyBcXG90aW1lcyBNJyJdLFswLDIsIk0iXSxbMiwyLCJNJyJdLFszLDEsIlxcZXF1aXYiXSxbNCwwLCJNIFxcb3RpbWVzIE0iXSxbNiwwLCJNJyBcXG90aW1lcyBNJyJdLFs0LDIsIk0iXSxbNiwyLCJNJyJdLFswLDEsImcgXFxvdGltZXMgZyIsMSx7ImN1cnZlIjoyfV0sWzAsMiwiXFxtdSIsMl0sWzIsMywiZyIsMl0sWzEsMywiXFxtdSciXSxbMCwxLCJmIFxcb3RpbWVzIGYiLDAseyJjdXJ2ZSI6LTV9XSxbMSwyLCJcXGNoaV97Z30iLDEseyJzaG9ydGVuIjp7InNvdXJjZSI6MjAsInRhcmdldCI6MjB9LCJsZXZlbCI6Mn1dLFs1LDYsImYgXFxvdGltZXMgZiJdLFs1LDcsIlxcbXUiLDJdLFs3LDgsImYiLDEseyJjdXJ2ZSI6LTJ9XSxbNiw4LCJcXG11JyJdLFs2LDcsIlxcY2hpX3tmfSIsMSx7InNob3J0ZW4iOnsic291cmNlIjoyMCwidGFyZ2V0IjoyMH0sImxldmVsIjoyfV0sWzcsOCwiZyIsMSx7ImN1cnZlIjo1fV0sWzEzLDksInIgXFxvdGltZXMgciIsMCx7InNob3J0ZW4iOnsic291cmNlIjoyMCwidGFyZ2V0IjoyMH19XSxbMTcsMjAsInIiLDAseyJzaG9ydGVuIjp7InNvdXJjZSI6MjAsInRhcmdldCI6MjB9fV1d
		\[\begin{tikzcd}[ampersand replacement=\&]
			{M \otimes M} \&\& {M' \otimes M'} \&\& {M \otimes M} \&\& {M' \otimes M'} \\
			\&\&\& \equiv \\
			M \&\& {M'} \&\& M \&\& {M'}
			\arrow[""{name=0, anchor=center, inner sep=0}, "{g \otimes g}"{description}, curve={height=12pt}, from=1-1, to=1-3]
			\arrow[""{name=1, anchor=center, inner sep=0}, "{f \otimes f}", curve={height=-30pt}, from=1-1, to=1-3]
			\arrow["\mu"', from=1-1, to=3-1]
			\arrow["{\chi_{g}}"{description}, between={0.2}{0.8}, Rightarrow, nfold, from=1-3, to=3-1]
			\arrow["{\mu'}", from=1-3, to=3-3]
			\arrow["{f \otimes f}", from=1-5, to=1-7]
			\arrow["\mu"', from=1-5, to=3-5]
			\arrow["{\chi_{f}}"{description}, between={0.2}{0.8}, Rightarrow, nfold, from=1-7, to=3-5]
			\arrow["{\mu'}", from=1-7, to=3-7]
			\arrow["g"', from=3-1, to=3-3]
			\arrow[""{name=2, anchor=center, inner sep=0}, "f"{description}, curve={height=-12pt}, from=3-5, to=3-7]
			\arrow[""{name=3, anchor=center, inner sep=0}, "g"{description}, curve={height=30pt}, from=3-5, to=3-7]
			\arrow["{r \otimes r}", between={0.2}{0.8}, Rightarrow, nfold, from=1, to=0]
			\arrow["r", between={0.2}{0.8}, Rightarrow, nfold, from=2, to=3]
		\end{tikzcd}\]
		% https://q.uiver.app/#q=WzAsNyxbMCwwLCJJIl0sWzAsMiwiTSJdLFsyLDIsIk0nIl0sWzMsMSwiXFxlcXVpdiJdLFs0LDAsIkkiXSxbNCwyLCJNIl0sWzYsMiwiTSciXSxbMCwxLCJlIiwyXSxbMCwyLCJlJyJdLFsxLDIsImYiLDEseyJjdXJ2ZSI6LTF9XSxbMSwyLCJnIiwyLHsiY3VydmUiOjV9XSxbNCw1LCJlIiwyXSxbNCw2LCJlJyJdLFs1LDYsImciLDJdLFs4LDEsIlxcaW90YV97Zn0iLDAseyJzaG9ydGVuIjp7InNvdXJjZSI6MjAsInRhcmdldCI6MjB9fV0sWzksMTAsInIiLDIseyJzaG9ydGVuIjp7InNvdXJjZSI6MjAsInRhcmdldCI6MjB9fV0sWzEyLDUsIlxcaW90YV97Z30iLDAseyJzaG9ydGVuIjp7InNvdXJjZSI6MjAsInRhcmdldCI6MjB9fV1d
		\[\begin{tikzcd}[ampersand replacement=\&]
			I \&\&\&\& I \&\& \\
			\&\&\& \equiv \\
			M \&\& {M'} \&\& M \&\& {M'}
			\arrow["e"', from=1-1, to=3-1]
			\arrow[""{name=0, anchor=center, inner sep=0}, "{e'}", from=1-1, to=3-3]
			\arrow["e"', from=1-5, to=3-5]
			\arrow[""{name=1, anchor=center, inner sep=0}, "{e'}", from=1-5, to=3-7]
			\arrow[""{name=2, anchor=center, inner sep=0}, "f"{description}, curve={height=-6pt}, from=3-1, to=3-3]
			\arrow[""{name=3, anchor=center, inner sep=0}, "g"', curve={height=30pt}, from=3-1, to=3-3]
			\arrow["g"', from=3-5, to=3-7]
			\arrow["{\iota_{f}}", between={0.2}{0.8}, Rightarrow, nfold, from=0, to=3-1]
			\arrow["{\iota_{g}}", between={0.2}{0.8}, Rightarrow, nfold, from=1, to=3-5]
			\arrow["r"', between={0.2}{0.8}, Rightarrow, nfold, from=2, to=3]
		\end{tikzcd}\cdpunct[8pt]{.}\]
	\end{enumerate}
\end{definition}

\begin{lemma}
	The data of pseudomonoids and their morphisms in $\cV$ assemble into a bicategory of pseudomonoids called $\nMon{1}(\cV)$. 
\end{lemma}

\begin{proof}
	This lemma primarily follows from \cite[Sec.~3]{gurski2014itericon} by considering one object version of iconic weak enrichment over a monoidal bicategory. 
\end{proof}

\begin{example}
	As expected, if $\cV$ is the monoidal bicategory $\nCat{1}$, then the pseudomonoid objects and their morphisms are usual monoidal categories and their lax monoidal functors.
\end{example}

Now, we shall describe braided pseudomonoid objects and their morphisms. For this, we require a braided monoidal bicategory. The details are ultimately irrelevant for us but can be found carefully constructed in \cite{gurski2011monbicat}. In particular, $\cV$ is a braided monoidal bicategory if it is a monoidal bicategory equipped with a family of adjoint equivalences $b_{x,y} : X \otimes Y \leftrightarrow Y \otimes X : b_{x,y}^{-1}$ that behaves as the braiding.

\begin{definition}
	A braided pseudomonoid in a braided monoidal bicategory $\cV$ is a pseudomonoid object $M$ together with the additional data of
	\begin{enumerate}
		\item a braiding $2$-cell $b$
		% https://q.uiver.app/#q=WzAsMyxbMCwwLCJNIFxcb3RpbWVzIE0iXSxbMiwwLCJNIFxcb3RpbWVzIE0iXSxbMSwyLCJNIl0sWzAsMSwiYl97TSxNfSJdLFswLDIsIlxcbXUiLDJdLFsxLDIsIlxcbXUiXSxbNCwxLCJiIiwyLHsic2hvcnRlbiI6eyJzb3VyY2UiOjIwLCJ0YXJnZXQiOjIwfX1dXQ==
		\[\begin{tikzcd}[ampersand replacement=\&]
			{M \otimes M} \&\& {M \otimes M} \\
			\\
			\& M
			\arrow["{b_{M,M}}", from=1-1, to=1-3]
			\arrow[""{name=0, anchor=center, inner sep=0}, "\mu"', from=1-1, to=3-2]
			\arrow["\mu", from=1-3, to=3-2]
			\arrow["b"', between={0.2}{0.8}, Rightarrow, nfold, from=0, to=1-3]
		\end{tikzcd}\cdpunct[8pt]{,}\]
		\item such that the two hexagon axioms (\cref{app:hexagon}) hold.
	\end{enumerate}
\end{definition}


\begin{definition}\leavevmode
	\begin{enumerate}
		\item A $1$-morphism $(f,\chi,i) : (M,b) \rightarrow (M',b')$ of braided pseudomonoid is a $1$-morphism of pseudomonoid together with the following addtional axiom:
		% https://q.uiver.app/#q=WzAsMTIsWzAsMSwiTSBcXG90aW1lcyBNIl0sWzEsMywiTSdcXG90aW1lcyBNJyJdLFswLDQsIk0iXSxbMiw0LCJNJyJdLFsyLDEsIk0nIFxcb3RpbWVzIE0nIl0sWzEsMCwiTSBcXG90aW1lcyBNIl0sWzMsMiwiXFxlcXVpdiJdLFs0LDEsIk0gXFxvdGltZXMgTSJdLFs1LDAsIk0gXFxvdGltZXMgTSJdLFs0LDQsIk0iXSxbNiwxLCJNJyBcXG90aW1lcyBNJyJdLFs2LDQsIk0nIl0sWzAsMSwiZiBcXG90aW1lcyBmIiwxXSxbMCwyLCJcXG11IiwyXSxbMiwzLCJmIiwyXSxbMSwzLCJcXG11JyIsMV0sWzEsMiwiXFxjaGkiLDIseyJzaG9ydGVuIjp7InNvdXJjZSI6MjAsInRhcmdldCI6MjB9LCJsZXZlbCI6Mn1dLFs0LDMsIlxcbXUnIl0sWzEsNCwiYl97TScsTSd9IiwxXSxbMCw1LCJiX3tNLE19Il0sWzUsNCwiZiBcXG90aW1lcyBmIl0sWzcsOCwiYl97TSxNfSJdLFs3LDksIlxcbXUiLDJdLFs4LDksIlxcbXUiXSxbOCwxMCwiZiBcXG90aW1lcyBmIl0sWzEwLDExLCJcXG11JyJdLFs5LDExLCJmIiwyXSxbMTAsOSwiXFxjaGkiLDAseyJzaG9ydGVuIjp7InNvdXJjZSI6MjAsInRhcmdldCI6MjB9LCJsZXZlbCI6Mn1dLFs0LDE1LCJiJyIsMSx7InNob3J0ZW4iOnsic291cmNlIjoyMCwidGFyZ2V0IjoyMH19XSxbOCwyMiwiYiIsMSx7InNob3J0ZW4iOnsic291cmNlIjoyMCwidGFyZ2V0IjoyMH19XV0=
		\[\begin{tikzcd}[ampersand replacement=\&,column sep=tiny]
			\& {M \otimes M} \&\&\&\& {M \otimes M} \& \\
			{M \otimes M} \&\& {M' \otimes M'} \&\& {M \otimes M} \&\& {M' \otimes M'} \\
			\&\&\& \equiv \\
			\& {M'\otimes M'} \\
			M \&\& {M'} \&\& M \&\& {M'}
			\arrow["{f \otimes f}", from=1-2, to=2-3]
			\arrow["{f \otimes f}", from=1-6, to=2-7]
			\arrow["\mu", from=1-6, to=5-5]
			\arrow["{b_{M,M}}", from=2-1, to=1-2]
			\arrow["{f \otimes f}"{description}, from=2-1, to=4-2]
			\arrow["\mu"', from=2-1, to=5-1]
			\arrow["{\mu'}", from=2-3, to=5-3]
			\arrow["{b_{M,M}}", from=2-5, to=1-6]
			\arrow[""{name=0, anchor=center, inner sep=0}, "\mu"', from=2-5, to=5-5]
			\arrow["\chi", between={0.2}{0.8}, Rightarrow, nfold, from=2-7, to=5-5]
			\arrow["{\mu'}", from=2-7, to=5-7]
			\arrow["{b_{M',M'}}"{description}, from=4-2, to=2-3]
			\arrow["\chi"', between={0.2}{0.8}, Rightarrow, nfold, from=4-2, to=5-1]
			\arrow[""{name=1, anchor=center, inner sep=0}, "{\mu'}"{description}, from=4-2, to=5-3]
			\arrow["f"', from=5-1, to=5-3]
			\arrow["f"', from=5-5, to=5-7]
			\arrow["b"{description}, between={0.2}{0.8}, Rightarrow, nfold, from=1-6, to=0]
			\arrow["{b'}"{description}, between={0.2}{0.8}, Rightarrow, nfold, from=2-3, to=1]
		\end{tikzcd}\cdpunct[8pt]{.}\]
		
		\item A $2$-morphism of braided pseudomonoid is the same data as that of a pseudomonoid.
	\end{enumerate}
\end{definition}

\begin{lemma}
	There is a bicategory of braided pseudomonoids $\nMon{1b}(\cV)$.
\end{lemma}

\begin{proof}
	Using that the $2$-morphism of braided pseudomonoids are the same as that of pseudomonoid, we define the compositional data and the coherence constraint of the bicategory $\nMon{1b}(\cV)$ to be the same as that of $\nMon{1}(\cV)$. Thus, we only need to verify that the composite of $1$-morphism of braided pseudomonoid also satisfy the additional braiding axiom. Suppressing notations and labeling, this follows directly from the following pasting equalities:
	% https://q.uiver.app/#q=WzAsMzUsWzAsMywiTSJdLFsxLDMsIk4iXSxbMiwzLCJQIl0sWzAsMSwiTU0iXSxbMSwyLCJOTiJdLFsyLDEsIlBQIl0sWzMsMCwiUFAiXSxbMSwwLCJNTSJdLFsxLDQsIlxcdmVxdWl2Il0sWzAsOSwiTSJdLFsxLDksIk4iXSxbMiw5LCJQIl0sWzAsNywiTU0iXSxbMiw3LCJQUCJdLFswLDUsIk1NIl0sWzMsNiwiUFAiXSxbMSw3LCJOTiJdLFsxLDYsIk5OIl0sWzQsNywiXFxlcXVpdiJdLFs1LDcsIk1NIl0sWzUsOSwiTSJdLFs2LDcsIk5OIl0sWzYsOSwiTiJdLFs3LDYsIk5OIl0sWzksNiwiUFAiXSxbOCw5LCJQIl0sWzUsNSwiTU0iXSxbNyw0LCJcXHZlcXVpdiJdLFs1LDMsIk0iXSxbNSwxLCJNTSJdLFs2LDAsIk1NIl0sWzcsMSwiTk4iXSxbOSwxLCJQUCJdLFs2LDMsIk4iXSxbOCwzLCJQIl0sWzUsNl0sWzMsNF0sWzUsMl0sWzQsMV0sWzMsMF0sWzAsMV0sWzEsMl0sWzQsNV0sWzYsMl0sWzMsNV0sWzMsN10sWzcsNl0sWzEyLDE0XSxbMTQsMTddLFsxMywxNV0sWzE3LDE1XSxbMTQsMTVdLFsxMiwxNl0sWzE2LDEzXSxbMTYsMTddLFsxNSwxMV0sWzEzLDExXSxbMTYsMTBdLFsxMiw5XSxbOSwxMF0sWzEwLDExXSxbMjEsMjNdLFsyMywyMl0sWzIzLDI0XSxbMjQsMjVdLFsyMiwyNV0sWzIxLDIyXSxbMTksMjFdLFsxOSwyMF0sWzIwLDIyXSxbMTksMjZdLFsyNiwyM10sWzI2LDI0XSxbMzAsMjhdLFsyOSwyOF0sWzI5LDMwXSxbMzAsMzFdLFszMSwzMl0sWzMwLDMyXSxbMjgsMzNdLFszMSwzM10sWzMzLDM0XSxbMzIsMzRdXQ==
	\[\begin{tikzcd}[ampersand replacement=\&,column sep=small]
		\& MM \&\& PP \&\&\& MM \&\&\& \\
		MM \&\& PP \&\&\& MM \&\& NN \&\& PP \\
		\& NN \\
		M \& N \& P \&\&\& M \& N \&\& P \\
		\& \vequiv \&\&\&\&\&\& \vequiv \\
		MM \&\&\&\&\& MM \\
		\& NN \&\& PP \&\&\&\& NN \&\& PP \\
		MM \& NN \& PP \&\& \equiv \& MM \& NN \\
		\\
		M \& N \& P \&\&\& M \& N \&\& P
		\arrow[from=1-2, to=1-4]
		\arrow[from=1-4, to=4-3]
		\arrow[from=1-7, to=2-8]
		\arrow[from=1-7, to=2-10]
		\arrow[from=1-7, to=4-6]
		\arrow[from=2-1, to=1-2]
		\arrow[from=2-1, to=2-3]
		\arrow[from=2-1, to=3-2]
		\arrow[from=2-1, to=4-1]
		\arrow[from=2-3, to=1-4]
		\arrow[from=2-3, to=4-3]
		\arrow[from=2-6, to=1-7]
		\arrow[from=2-6, to=4-6]
		\arrow[from=2-8, to=2-10]
		\arrow[from=2-8, to=4-7]
		\arrow[from=2-10, to=4-9]
		\arrow[from=3-2, to=2-3]
		\arrow[from=3-2, to=4-2]
		\arrow[from=4-1, to=4-2]
		\arrow[from=4-2, to=4-3]
		\arrow[from=4-6, to=4-7]
		\arrow[from=4-7, to=4-9]
		\arrow[from=6-1, to=7-2]
		\arrow[from=6-1, to=7-4]
		\arrow[from=6-6, to=7-8]
		\arrow[from=6-6, to=7-10]
		\arrow[from=7-2, to=7-4]
		\arrow[from=7-4, to=10-3]
		\arrow[from=7-8, to=7-10]
		\arrow[from=7-8, to=10-7]
		\arrow[from=7-10, to=10-9]
		\arrow[from=8-1, to=6-1]
		\arrow[from=8-1, to=8-2]
		\arrow[from=8-1, to=10-1]
		\arrow[from=8-2, to=7-2]
		\arrow[from=8-2, to=8-3]
		\arrow[from=8-2, to=10-2]
		\arrow[from=8-3, to=7-4]
		\arrow[from=8-3, to=10-3]
		\arrow[from=8-6, to=6-6]
		\arrow[from=8-6, to=8-7]
		\arrow[from=8-6, to=10-6]
		\arrow[from=8-7, to=7-8]
		\arrow[from=8-7, to=10-7]
		\arrow[from=10-1, to=10-2]
		\arrow[from=10-2, to=10-3]
		\arrow[from=10-6, to=10-7]
		\arrow[from=10-7, to=10-9]
	\end{tikzcd}\cdpunct[8pt]{.}\]
\end{proof}

Finally, we can also describe symmetric pseudomonoid object inside a sylleptic (or symmetric) monoidal bicategory. This bicategory is properly defined in \cite{gurski2013symmonbicat} and is braided monoidal bicategory together with a $2$-cell $\upsilon_{X,Y} : b_{Y,X}b_{X,Y} \Rightarrow 1$ that essentially identifies $b_{X,Y}$ with $b^{-1}_{Y,X}$. 

\begin{definition}\leavevmode
	\begin{enumerate}
		\item A symmetric pseudomonoid is a braided pseudmonoid that satisfy the following additional axiom:
		% https://q.uiver.app/#q=WzAsOCxbMCwwLCJNIFxcb3RpbWVzIE0iXSxbMiwxLCJNIFxcb3RpbWVzIE0iXSxbNCwwLCJNIFxcb3RpbWVzIE0iXSxbMiwzLCJNIl0sWzUsMSwiXFxlcXVpdiJdLFs2LDAsIk0gXFxvdGltZXMgTSJdLFs4LDAsIk0gXFxvdGltZXMgTSJdLFs3LDMsIk0iXSxbMCwzLCJcXG11IiwyXSxbMCwxLCJiX3tNLE19IiwxXSxbMSwyLCJiX3tNLE19IiwxXSxbMiwzLCJcXG11Il0sWzEsM10sWzAsMiwiXFx2YXJlcHNpbG9uX3tNIFxcb3RpbWVzIE19Il0sWzUsNywiXFxtdSIsMl0sWzYsNywiXFxtdSJdLFs1LDYsIlxcdmFyZXBzaWxvbl97TSBcXG90aW1lcyBNfSJdLFs4LDEyLCJiIiwwLHsic2hvcnRlbiI6eyJzb3VyY2UiOjIwLCJ0YXJnZXQiOjIwfX1dLFsxMiwxMSwiYiIsMCx7InNob3J0ZW4iOnsic291cmNlIjoyMCwidGFyZ2V0IjoyMH19XSxbMSwxMywiXFx1cHNpbG9uX3tNLE19IiwwLHsic2hvcnRlbiI6eyJzb3VyY2UiOjIwLCJ0YXJnZXQiOjIwfX1dLFsxNCwxNSwiMSIsMix7InNob3J0ZW4iOnsic291cmNlIjoyMCwidGFyZ2V0IjoyMH19XV0=
		\[\begin{tikzcd}[ampersand replacement=\&,column sep=tiny]
			{M \otimes M} \&\&\&\& {M \otimes M} \&\& {M \otimes M} \&\& {M \otimes M} \\
			\&\& {M \otimes M} \&\&\& \equiv \\
			\\
			\&\& M \&\&\&\&\& M
			\arrow[""{name=0, anchor=center, inner sep=0}, "{\varepsilon_{M \otimes M}}", from=1-1, to=1-5]
			\arrow["{b_{M,M}}"{description}, from=1-1, to=2-3]
			\arrow[""{name=1, anchor=center, inner sep=0}, "\mu"', from=1-1, to=4-3]
			\arrow[""{name=2, anchor=center, inner sep=0}, "\mu", from=1-5, to=4-3]
			\arrow["{\varepsilon_{M \otimes M}}", from=1-7, to=1-9]
			\arrow[""{name=3, anchor=center, inner sep=0}, "\mu"', from=1-7, to=4-8]
			\arrow[""{name=4, anchor=center, inner sep=0}, "\mu", from=1-9, to=4-8]
			\arrow["{b_{M,M}}"{description}, from=2-3, to=1-5]
			\arrow[""{name=5, anchor=center, inner sep=0}, from=2-3, to=4-3]
			\arrow["b", between={0.2}{0.8}, Rightarrow, nfold, from=1, to=5]
			\arrow["1"', between={0.2}{0.8}, Rightarrow, nfold, from=3, to=4]
			\arrow["{\upsilon_{M,M}}", between={0.2}{0.8}, Rightarrow, nfold, from=2-3, to=0]
			\arrow["b", between={0.2}{0.8}, Rightarrow, nfold, from=5, to=2]
		\end{tikzcd}\cdpunct[8pt]{.}\]
		\item The $1$-morphisms and $2$-morphisms of symmetric pseudomonoid object has no additional data.
	\end{enumerate}
\end{definition}

\begin{lemma}
	There is a bicategory of symmetric pseudomonoids $\nMon{1s}(\cV)$, which is formed as a full sub-bicategory of braided pseudomonoids.
\end{lemma}

For the case of $\cV=\nCat{1}$, we know that $\nCat{1}$ is a symmetric monoidal $2$-category. Thus, as expected, $\nMon{1b}(\nCat{1}) = \BrMonCat_l$ and $\nMon{1s}(\nCat{1})= \SymMonCat_l$. Similarly, it is possible to consider colax $1$-morphisms in the constructions of $\nMon{1}$ and its variants. In this case, we obtain the corresponding $2$-category of monoidal categories with colax monoidal functors.

With all the relevant definitions, we can now define pseudomodules over pseudomonoid objects.

\begin{definition}
	A left pseudomodule object over a pseudomonoid object $(M,\mu,e, a_{M},\lambda_{M},\rho_{M})$ in a monoidal bicategory $\cV$ is given by the collection $(\cZ,\triangleleft,a_{\cZ},\lambda_{\cZ})$ such that:
	\begin{enumerate}
		\item $\cZ \in \cV$ is an object, $\triangleleft$ defines a left action $1$-cell $M \otimes \cZ \xrightarrow{\triangleleft} \cZ$ in $\cV$,
		\item the associator $a_{\cZ}$ and left unitor $\lambda_{\cZ}$ are invertible $2$-cells in $\cV$ which are defined using the following diagrams:
		% https://q.uiver.app/#q=WzAsOCxbMCwwLCJNIFxcb3RpbWVzIChNIFxcb3RpbWVzIFxcY1opIl0sWzAsMiwiKE0gXFxvdGltZXMgTSkgXFxvdGltZXMgXFxjWiJdLFswLDQsIk0gXFxvdGltZXMgXFxjWiJdLFsyLDQsIlxcY1oiXSxbMiwwLCJNIFxcb3RpbWVzIFxcY1oiXSxbNCw0LCJJIFxcb3RpbWVzIFxcY1oiXSxbNCwwLCJNIFxcb3RpbWVzIFxcY1oiXSxbNiwwLCJcXGNaIl0sWzAsNCwiMSBcXG90aW1lcyBcXHRyaWFuZ2xlbGVmdCJdLFs0LDMsIlxcdHJpYW5nbGVsZWZ0Il0sWzIsMywiXFx0cmlhbmdsZWxlZnQiLDJdLFsxLDIsIlxcbXUgXFxvdGltZXMgMSIsMl0sWzEsMCwiYV97XFxjVn0iXSxbNSw2LCJlIFxcb3RpbWVzIDEiXSxbNSw3LCJcXGxhbWJkYV97XFxjVn0iLDJdLFs2LDcsIlxcdHJpYW5nbGVsZWZ0Il0sWzIsNCwiYV97XFxjWn0iLDIseyJzaG9ydGVuIjp7InNvdXJjZSI6MjAsInRhcmdldCI6MjB9LCJsZXZlbCI6Mn1dLFs2LDE0LCJcXGxhbWJkYV97XFxjWn0iLDIseyJzaG9ydGVuIjp7InNvdXJjZSI6MjAsInRhcmdldCI6MjB9fV1d
		\[\begin{tikzcd}[ampersand replacement=\&]
			{M \otimes (M \otimes \cZ)} \&\& {M \otimes \cZ} \&\& {M \otimes \cZ} \&\& \cZ \\
			\\
			{(M \otimes M) \otimes \cZ} \\
			\\
			{M \otimes \cZ} \&\& \cZ \&\& {I \otimes \cZ}
			\arrow["{1 \otimes \triangleleft}", from=1-1, to=1-3]
			\arrow["\triangleleft", from=1-3, to=5-3]
			\arrow["\triangleleft", from=1-5, to=1-7]
			\arrow["{a_{\cV}}", from=3-1, to=1-1]
			\arrow["{\mu \otimes 1}"', from=3-1, to=5-1]
			\arrow["{a_{\cZ}}"', between={0.2}{0.8}, Rightarrow, nfold, from=5-1, to=1-3]
			\arrow["\triangleleft"', from=5-1, to=5-3]
			\arrow["{e \otimes 1}", from=5-5, to=1-5]
			\arrow[""{name=0, anchor=center, inner sep=0}, "{\lambda_{\cV}}"', from=5-5, to=1-7]
			\arrow["{\lambda_{\cZ}}"', between={0.2}{0.8}, Rightarrow, nfold, from=1-5, to=0]
		\end{tikzcd}\cdpunct[8pt]{,}\]
		\item the pentagon axiom (\cref{app:pentagon}) and the triangle axiom (\cref{app:triangle}) hold with replacing the corresponding copy of $M$ with $\cZ$. 
	\end{enumerate}
\end{definition}

\begin{definition}
	A right pseudomodule object over a pseudomonoid object $(M,\mu,e, a_{M},\lambda_{M},\rho_{M})$ in a monoidal bicatgory $\cV$ is given by the collection $(\cZ,\triangleright,a_{\cZ},\rho_{\cZ})$ such that:
	\begin{enumerate}
		\item $\cZ \in \cV$ is an object, $\triangleright$ defines a right action $1$-cell $\cZ \otimes M \xrightarrow{\triangleright} \cZ$ in $\cV$,
		\item the associator $a_{\cZ}$ and right unitor $\rho_{\cZ}$ are invertible $2$-morphisms in $\cV$ which are defined using the following diagrams:
		% https://q.uiver.app/#q=WzAsOCxbMCwwLCIgKFxcY1ogXFxvdGltZXMgTSkgXFxvdGltZXMgTSJdLFswLDIsIlxcY1ogXFxvdGltZXMgKE0gXFxvdGltZXMgTSkiXSxbMCw0LCJcXGNaIFxcb3RpbWVzIE0iXSxbMiw0LCJcXGNaIl0sWzIsMCwiXFxjWiBcXG90aW1lcyBNIl0sWzQsNCwiXFxjWiBcXG90aW1lcyBJIl0sWzQsMCwiXFxjWiBcXG90aW1lcyBNIl0sWzYsMCwiXFxjWiJdLFswLDQsIlxcdHJpYW5nbGVyaWdodCBcXG90aW1lcyAxIl0sWzQsMywiXFx0cmlhbmdsZXJpZ2h0Il0sWzIsMywiXFx0cmlhbmdsZXJpZ2h0IiwyXSxbMSwyLCIxIFxcb3RpbWVzIFxcbXUiLDJdLFswLDEsImFfe1xcY1Z9IiwyXSxbNSw2LCIxIFxcb3RpbWVzIGUiXSxbNSw3LCJcXHJob197XFxjVn0iLDJdLFs2LDcsIlxcdHJpYW5nbGVyaWdodCJdLFs0LDIsImFfe1xcY1p9IiwwLHsic2hvcnRlbiI6eyJzb3VyY2UiOjIwLCJ0YXJnZXQiOjIwfSwibGV2ZWwiOjJ9XSxbNiwxNCwiXFxyaG9fe1xcY1p9IiwyLHsic2hvcnRlbiI6eyJzb3VyY2UiOjIwLCJ0YXJnZXQiOjIwfX1dXQ==
		\[\begin{tikzcd}[ampersand replacement=\&]
			{ (\cZ \otimes M) \otimes M} \&\& {\cZ \otimes M} \&\& {\cZ \otimes M} \&\& \cZ \\
			\\
			{\cZ \otimes (M \otimes M)} \\
			\\
			{\cZ \otimes M} \&\& \cZ \&\& {\cZ \otimes I}
			\arrow["{\triangleright \otimes 1}", from=1-1, to=1-3]
			\arrow["{a_{\cV}}"', from=1-1, to=3-1]
			\arrow["{a_{\cZ}}", between={0.2}{0.8}, Rightarrow, nfold, from=1-3, to=5-1]
			\arrow["\triangleright", from=1-3, to=5-3]
			\arrow["\triangleright", from=1-5, to=1-7]
			\arrow["{1 \otimes \mu}"', from=3-1, to=5-1]
			\arrow["\triangleright"', from=5-1, to=5-3]
			\arrow["{1 \otimes e}", from=5-5, to=1-5]
			\arrow[""{name=0, anchor=center, inner sep=0}, "{\rho_{\cV}}"', from=5-5, to=1-7]
			\arrow["{\rho_{\cZ}}"', between={0.2}{0.8}, Rightarrow, nfold, from=1-5, to=0]
		\end{tikzcd}\cdpunct[8pt]{,}\]
		\item the pentagon axiom (\cref{app:pentagon}) and the triangle axiom (\cref{app:triangle}) hold with replacing the corresponding copy of $M$ with $\cZ$. 
	\end{enumerate}
\end{definition}

\begin{definition}
	A bipseudomodule $\cZ \in \cV$ between two pseudomonoids $(M,\mu_{M},e_{M}, a_{M},\lambda_{M},\rho_{M})$ and $(N,\mu_{N},e_{N}, a_{N},\lambda_{N},\rho_{N})$ is given by 
	\begin{enumerate}
		\item a left pseudomodule $(\cZ,\triangleleft,a_{(\cZ,M)},\lambda_{\cZ})$ over $M$,
		\item a right pseudomodule $(\cZ,\triangleright,a_{(\cZ,N)},\rho_{\cZ})$ over $N$,
		\item an invertible balancing $2$-cell $l$ in $\cV$ which is defined using the following diagram:
		% https://q.uiver.app/#q=WzAsNSxbMCwwLCIgKE0gXFxvdGltZXMgXFxjWikgXFxvdGltZXMgTiJdLFs0LDAsIk0gXFxvdGltZXMgKFxcY1ogXFxvdGltZXMgTikiXSxbNCwyLCJNIFxcb3RpbWVzIFxcY1oiXSxbMiw0LCJcXGNaIl0sWzAsMiwiXFxjWiBcXG90aW1lcyBOIl0sWzAsNCwiXFx0cmlhbmdsZWxlZnRfe1xcY1p9XFxvdGltZXMgMSIsMl0sWzQsMywiXFx0cmlhbmdsZXJpZ2h0X3tcXGNafSIsMl0sWzIsMywiXFx0cmlhbmdsZWxlZnRfe1xcY1p9Il0sWzEsMiwiMSBcXG90aW1lcyBcXHRyaWFuZ2xlcmlnaHRfe1xcY1p9Il0sWzAsMSwiYV97XFxjVn0iXSxbNCwyLCJsIiwwLHsic2hvcnRlbiI6eyJzb3VyY2UiOjIwLCJ0YXJnZXQiOjIwfSwibGV2ZWwiOjJ9XV0=
		\[\begin{tikzcd}[ampersand replacement=\&,cramped]
			{ (M \otimes \cZ) \otimes N} \&\&\&\& {M \otimes (\cZ \otimes N)} \\
			\\
			{\cZ \otimes N} \&\&\&\& {M \otimes \cZ} \\
			\\
			\&\& \cZ
			\arrow["{a_{\cV}}", from=1-1, to=1-5]
			\arrow["{\triangleleft_{\cZ}\otimes 1}"', from=1-1, to=3-1]
			\arrow["{1 \otimes \triangleright_{\cZ}}", from=1-5, to=3-5]
			\arrow["l", between={0.2}{0.8}, Rightarrow, nfold, from=3-1, to=3-5]
			\arrow["{\triangleright_{\cZ}}"', from=3-1, to=5-3]
			\arrow["{\triangleleft_{\cZ}}", from=3-5, to=5-3]
		\end{tikzcd}\cdpunct[8pt]{,}\]
		\item two instances of the pentagon axiom (\cref{app:pentagon}) hold, one for $(M,\cZ,N,N)$ and the other for $(M,M,\cZ,N)$.  
	\end{enumerate}
\end{definition}

All the axioms for left, right, and bi-pseudomodule are motivated by viewing monoidal categories and module categories as parts of a two objects bicategory $K$. Furthermore, in the case of $\cV=\nCat{1}$, these provide the usual definitions.


